\section{Horizontal pluriharmonic currents}
\label{sec:horizontal-pluriharmonic-currents}

A current is \emph{pluriharmonic} if its $\ddc$ vanishes. A vanishing vertical trace removes the vertical and mixed coefficients of a positive $(1,1)$-current. For a closed current the remaining coefficients are constant along the fibers. For a pluriharmonic current they are harmonic; positivity gives the same descent when the fiber is the whole complex line or a compact connected curve.

\begin{lemma}\label{lem:horizontal-pluriharmonic-descent}\index{Positive current!horizontal descent}\index{Pluriharmonic current!horizontal descent}
Let $Y$ be a complex manifold of dimension $\ell\geq1$, let $\eta$ be a Hermitian form on $Y$, and let $\pi:Z\to Y$ be a holomorphic fiber bundle with fiber either $\mathbb C$ or a compact connected Riemann surface. If $S$ is a positive real $(1,1)$-current on $Z$ with $\ddc S=0$ and $S\wedge\pi^*\eta^\ell=0$, then there is a unique positive $\ddc$-closed $(1,1)$-current $R$ on $Y$ such that $S=\pi^*R$. If $S$ is closed, then $R$ is closed.
\end{lemma}
\begin{proof}
Fix a bundle trivialization $\pi^{-1}(U)=U\times F$, with base coordinates $w=(w_1,\ldots,w_\ell)$ and a local fiber coordinate $z$. The coefficient measures of $S$ have a positive Hermitian matrix of densities relative to their trace measure. The identity $S\wedge\pi^*\eta^\ell=0$ makes the vertical diagonal entry zero. The matrix inequality $|s_{\alpha\bar z}|^2\leq s_{\alpha\bar\alpha}s_{z\bar z}$ then removes all mixed entries. Thus
\[
 S=\frac{\mathrm{i}}{2\pi}\sum_{\alpha,\beta=1}^\ell s_{\alpha\bar\beta}(w,z)\,\mathrm{d}w_\alpha\wedge\mathrm{d}\bar w_\beta,
\]
with distribution coefficients. In $\partial\bar\partial S=0$, the coefficient containing both vertical differentials and the horizontal pair $(\alpha,\beta)$ gives $\partial_z\partial_{\bar z}s_{\alpha\bar\beta}=0$.

Fix $v\in\mathbb C^\ell$ and a nonnegative smooth compactly supported function $\xi$ on $U$. Pair the positive coefficient distribution $\sum_{\alpha,\beta}s_{\alpha\bar\beta}v_\alpha\bar v_\beta$ in the base variables against $\xi$ times the Euclidean volume form. This defines a nonnegative harmonic distribution on the whole fiber $F$. Its definition is independent of the fiber coordinate, since the remaining differentials are horizontal. Weyl's lemma makes it a nonnegative harmonic function. Such a function is constant on $\mathbb C$ by the harmonic Liouville theorem, and on a compact connected curve by the maximum principle.

Polarization in $v$ and linearity in the base test function show that every coefficient $s_{\alpha\bar\beta}$ is pulled back from a distribution $r_{\alpha\bar\beta}$ on $U$. To identify this distribution, pair against any nonnegative smooth fiber density of total integral one. This also proves positivity of the resulting current $R_U$ and injectivity of pullback. Consequently $\ddc S=\pi^*(\ddc R_U)=0$ implies $\ddc R_U=0$. On overlaps, injectivity makes the currents $R_U$ agree; they glue to $R$, and injectivity proves uniqueness. If $\mathrm{d}S=0$, the same argument gives $\mathrm{d}R=0$.
\end{proof}

The whole-fiber condition is essential. On the bidisc $\Delta_z\times\Delta_w$, the current
\[
 S=(2+\operatorname{Re}z)\frac{\mathrm{i}}{2\pi}\,\mathrm{d}w\wedge\mathrm{d}\bar w
\]
is positive and $\ddc$-closed, and its vertical trace vanishes. Its horizontal coefficient depends on $z$, so it does not descend under projection to $\Delta_w$.
